\subsection{Finite coverings}

A covering is said to be locally finite if the cardinalities of all its fibers are finite. A covering is said to be finite if the set of cardinalities of its fibers is bounded above by a finite cardinal.

THEOREM 1. --- Let E, B be topological spaces and \(p\colon\mathrm{E}\to\mathrm{B}\) a map. The following conditions are equivalent:

\begin{enumerate}
\def\labelenumi{(\roman{enumi})}
\item
  The space E, equipped with the map p, is a locally finite covering of B;
\item
  The map p is étale, proper, and separated;
\item
  The map p is continuous, open, and separated, its fibers are finite, and the numerical function \(b \mapsto \text{Card}(\overline{p}^{1}(b))\) is upper semicontinuous on B (TG, IV, p. 28).
\end{enumerate}

We will use the following lemma in the proof:

Lemma. — Let X and Y be topological spaces. For a map \(f: X \to Y\) to be closed, it is necessary and sufficient that for every point y of Y and every neighborhood W of the fiber \(f^{-1}(y)\) there exists a neighborhood V of y such that W contains \(f^{-1}(V)\) .

In this statement, one may consider only the neighborhoods W of \(f^{-1}(y)\) which are open. Denoting by F the complement of W in X, one can then reformulate the statement as follows: for a map \(f: X \to Y\) to be closed, it is necessary and sufficient that, for every closed subset F of X, every point y of Y that does not belong to \(f(F)\) has a neighborhood V disjoint from \(f(F)\). Now this assertion follows immediately from the definition of a closed map (TG, I, p. 30, def. 1).

Let us now prove Theorem 1. Each of the three conditions implies that the map p is continuous, open, and separated, and also that the fibers of p are finite. This is clear under hypotheses (i) and (iii); under hypothesis (ii), the fibers of p are discrete (I, p. 29, remark 2) and quasi-compact (TG, I, p. 75, th. 1), hence finite (TG, I, p. 60, example 1).

(i)⇒(ii): it suffices to show that p is proper and, for this, that for every open set U over which the covering (E,p) is trivializable, the map \(p_{\mathrm{U}}\colon \overline{p}^{1}(\mathrm{U})\to\mathrm{U}\) is proper (TG, I, p. 72, prop. 3).

As the fibers of \(p_{U}\) are finite, this last assertion follows from Corollary 5 of TG, I, p. 77.

(ii)⇒(iii): let b be a point of B and, for every \(x \in \overline{p}^{1}(b)\), let \(W_{x}\) be an open neighborhood of x in E such that \(p \mid W_{x}\) is injective. The set \(W = \bigcup_{x \in \overline{p}^{1}(b)} W_{x}\) is an open neighborhood of \(\overline{p}^{1}(b)\). Since the map p is closed, it follows from the lemma above that there exists an open neighborhood U of b such that \(\overline{p}^{1}(U) \subset W\). For all \(a \in U\), we have \(\overline{p}^{1}(a) \subset W\); since the restriction of p to each \(W_{x}\) is injective, it follows that \(\text{Card}(\overline{p}^{1}(a)) \leqslant \text{Card}(\overline{p}^{1}(b))\), which proves the upper semicontinuity of the map \(a \mapsto \text{Card}(\overline{p}^{1}(a))\).

(iii)⇒(i): let b be a point of B. Since the fiber \(E_{b} = \overline{p}^{1}(b)\) is finite and the map p is separated, one can choose, for every \(x \in E_{b}\), an open neighborhood \(V_{x}^{\prime}\) of x such that the \(V_{x}^{\prime}\) are pairwise disjoint (I, p. 26, Remark 4). Since the map p is open and the set \(E_{b}\) finite, the set \(U^{\prime} = \bigcap_{x \in E_{b}} p(V_{x}^{\prime})\) is an open neighborhood of b in B. Let U be an open neighborhood of b in B, contained in \(U^{\prime}\) and such that for all \(a \in U\), \(\text{Card}(E_{a}) \leqslant \text{Card}(E_{b})\). For all \(x \in E_{b}\), set \(V_{x} = V_{x}^{\prime} \cap \overline{p}^{1}(U)\). Let a be a point of U; the sets \(E_{a} \cap V_{x}\), for \(x \in E_{b}\), are nonempty and pairwise disjoint. These sets therefore each contain a unique element and form a partition of \(E_{a}\). This proves that, for every \(x \in E_{b}\), the map \(p | V_{x}\) is injective and that we have \(\overline{p}^{1}(U) = \bigcup_{x \in E_{b}} V_{x}\). Since the map p is open, it induces a homeomorphism from \(V_{x}\) onto U, and consequently, (E, p) is a covering of B (I, p. 70, cor. 1 of prop. 1).

Remark. --- A continuous, separated, open, proper map with finite fibers is not necessarily a covering. This is the case for the map \(x \mapsto x^{2}\) from R to \([0; +\infty[\) .

COROLLARY 1. --- Let E and B be topological spaces and let \(p\colon\mathrm{E}\to\mathrm{B}\) a proper and separated map. The set U of points b in B such that p is étale at every point of the fiber \(\mathrm{E}_{b}\) is open in B. The U-space induced by (E,p) over U is a locally finite covering.

The set F of points of E at which the map p is not étale is closed in E (I, p. 29, remark 1). Its image \(p(\mathrm{F})\) is closed because the map p is proper; the complement U of \(p(\mathrm{F})\) is therefore open. The map \(p_{U}\) is separated (I, p. 27, prop. 4) and proper (TG, I, p. 72, prop. 3). By construction, the map \(p_{\mathrm{U}}\) is étale; it therefore satisfies condition (ii) of Theorem 1.

COROLLARY 2. — Let B be a separated topological space and let E be a B-space. Suppose that E is compact and that its projection p: E → B is étale. Then E is a finite covering of B.

The map \(p\) is separated (I, p. 26, remark 2) and proper (TG, I, p. 76, cor. 2). By Corollary 1, E is therefore a locally finite covering of B. Since E is compact, \(p(\mathrm{E})\) is a quasi-compact subset of B; the map from \(p(\mathrm{E})\) into \(\mathbb{N}\) given by \(b \mapsto \operatorname{Card} \overline{p}^1(b)\) is locally constant, and is therefore bounded above (TG, IV, p. 30, corollary).

COROLLARY 3. --- Let B be a topological space, let (E,p) and (E',p') be B-spaces, and let f: E' → E be a B-morphism. Suppose that E' is a locally finite covering of B.

\begin{enumerate}
\def\labelenumi{\alph{enumi})}
\item
  If the map p is étale and separated, the E-space (E', f) is a covering.
\item
  If the map f is surjective and makes the space E' a covering of E, then E is a covering of B.
\end{enumerate}

Under the assumption \(a\), the maps \(p\) and \(p'\) are étale and separated, and the map \(p'\) is proper (Theorem 1). The map \(f\) is therefore étale (I, p. 30, cor. 1 of prop. 6), separated (I, p. 25, prop. 2 b)) and proper (I, p. 28, prop. 5). By Theorem 1, \((\mathrm{E}', f)\) is a covering of E.

Suppose now that \(f\) is surjective and suppose that \((\mathrm{E}', f)\) is a covering of E. It is then locally finite, so the map \(f\) is proper and surjective (th. 1). By I, p. 25, prop. 2, c), the map \(p\) is separated, since \(p' = p \circ f\) and \(p'\) is separated. The map \(p\) is étale (I, p. 29, prop. 6, d)), proper (TG, I, p. 73, prop. 5, b)), and its fibers are finite. By th. 1, the B-space \((\mathrm{E}, p)\) is then a covering of B.

COROLLARY 4. --- Let

\[
\begin{array}{c} \mathrm{E} ^ {\prime} \xrightarrow {f ^ {\prime}} \mathrm{E} \\ \Big \downarrow p ^ {\prime} \\ \mathrm{B} ^ {\prime} \xrightarrow {f} \mathrm{B} \end{array}
\]

a Cartesian square. Suppose that \((\mathrm{E}^{\prime}, p^{\prime})\) is a locally finite covering of \(B^{\prime}\) and that the map f is universally strict and surjective. Then \((\mathrm{E}, p)\) is a locally finite covering of B.

Indeed, the map p is étale (I, p. 31, prop. 8), proper (I, p. 21, prop. 11), and separated (I, p. 27, prop. 4).

COROLLARY 5. --- Let B be a connected topological space and let (E, p) be a finite covering of B. The space E has only finitely many connected components, and each of them is a covering of B.

The result is true if B is empty; we henceforth assume it to be nonempty. If X is a clopen subset of E, the restriction of p to X is a separated map (I, p. 25, prop. 2, a) and I, p. 26, remark 1), étale and proper (TG, I, p. 74, corollary 1); by Theorem 1, (X, p) is a locally finite covering of B. Since B is connected, this covering has finite degree. If X is distinct from E, there exists a point \(b \in B\) such that \(X_{b} \subsetneq E_{b}\), whence \([X:B] < [E:B]\). Since B is connected, it follows that every decreasing sequence of closed-open sets in E is stationary.

Let \(x \in E\); there then exists a smallest open and closed subset X of E containing x (E, III, p. 51, prop. 6). Such a set X is connected; it is therefore the connected component of x in E. This shows that the connected components of E are open and closed and that each of them is a covering of B. Since B is connected, each connected component of E meets every fiber of p; hence the connected components of E are finite in number.

PROPOSITION 5. — Let B be a topological space, let E and E′ be B-spaces, and let f: E′ → E be a B-morphism. Suppose that E is a locally finite covering and that the space E′, equipped with the map f, is a locally trivial E-fibered space (resp. a covering). Then E′ is a locally trivial B-fibered space (resp. a covering).

Denote by \(p\) and \(p'\) the respective projections of the B-spaces E and E'. Let \(b\) be a point of B. There exists an open neighborhood U of \(b\) in B and a finite partition \((\mathrm{V}_i)_{i\in \mathrm{I}}\) of \(\overline{p}^1 (\mathrm{U})\) formed by open sets of E such that, for every \(i\in \mathrm{I}\), the map \(p\) induces a homeomorphism of \(\mathrm{V}_i\) onto U. Set \(\mathrm{V}_i' = f^{-1}(\mathrm{V}_i)\). The sets \(\mathrm{V}_i'\) are open in E', form a partition of \(\overline{p'}^{1}(\mathrm{U})\), and the map from \(\mathrm{V}_i'\) to U induced by \(p'\) by passing to subspaces makes \(V_i'\) a locally trivial U-fibered space. It follows that the space \(\overline{p'}^{1}(\mathrm{U})\), equipped with the map \(p_U': \overline{p'}^{1}(\mathrm{U}) \to U\), is a locally trivial U-fibered space (I, p. 69, remark 2). If \((E', f)\) is a covering of E, the fibers of \(p_U'\) are discrete because their intersections with each of the open sets \(V_i'\) are discrete. This completes the proof.

Note. --- If (E, p) is a covering of B and (E', f) is a covering of E, the map \(p \circ f\) is étale (I, p. 29, prop. 6, a)) and separated (I, p. 25, prop. 2, a)). But it can happen, see Exercise 5, b) of I, p. 146, that the space E′, equipped with the map \(p \circ f\), is not a covering of B. See, however, IV, p. 342, prop. 3.

